\section{Choice after assignment}
A participant receives a role before choosing an action. Physical opponent
states belong to a finite set $T$; disclosed correspondences
$\theta_i:T\to A_{-i}$ fix their meaning in each candidate role. Opponents
follow a known law $\mu$ independently of the participant's choice. Receipts
$r_i(a,\theta_i(t))$ are measured in a common monetary scale, and the
participant is risk neutral. These are substantive calibration assumptions.
Put $x_{ia}(t)=r_i(a,\theta_i(t))$. The role value is
\[
 V_i(\mu)=\max_a\mu\cdot x_{ia}.
\]
A role lottery $\lambda$ is valued at $\sum_i\lambda_i V_i(\mu)$.
Indifference between all such lotteries is therefore exactly equality of the
role values. This definition makes no reference to automorphisms.

\begin{theorem}[Opportunity criterion; P7-A1]
Indifference holds for every belief on $T$ if and only if the sets
\[
 D_i=\operatorname{conv}\{x_{ia}:a\in A_i\}+\R_{-}^{T}
\]
are equal. It suffices to test all full-support beliefs.
\end{theorem}
\begin{proof}
The support function of $D_i$ at a probability vector is $V_i$.
If two closed downward polyhedra differ, separation supplies a separating
normal. Downward closure forces every coordinate of that normal to be
nonnegative. Normalize its sum to one to obtain a distinguishing belief.
Continuity extends equality on full-support beliefs to the whole simplex.
\end{proof}
This is an application of standard convex separation, not a new convexity
theorem. If every action is uniquely optimal at some full-support belief,
equality of these sets forces equality of payoff-row sets: unique exposed
maximizers of a common polyhedron coincide. Inessential actions need not be
identified. For example the menus
$\{(1,0),(0,1),(1/2,1/2)\}$ and
$\{(1,0),(0,1),(1/4,1/4)\}$ have equal values at every belief.
Neither added row is dominated by a single original row.

Even essential actions do not imply a simultaneous player automorphism.
For three binary players let $r=(1,0,0)$ and
$u_i(s)=\mathbf1\{s_i=s_{r(i)}\}$. In each role label the relevant opponent
first. All focal tables and optimized values coincide, and both actions are
essential. But the directed dependency graph has distinct indegrees $2,1,0$.
Every player permutation in an exact, cardinal or ordinal automorphism must
fix that graph, hence is the identity. Focal opportunity equality is weaker
than simultaneous typed interchangeability.

\paragraph{Information restrictions.}
For exchangeable beliefs, replace each payoff row by its averages on opponent
permutation orbits and apply the same theorem on the orbit simplex. Optimizing
before comparing expectations changes the relation to IID tests. On states
$00,01,10,11$, compare the menus
\[
 M_0=\{(2,0,0,0),(0,0,0,2),0\},\qquad
 M_1=\{(2,0,0,0),(0,0,0,2),(0,1,1,0)\}.
\]
Every common IID binary law gives both menus value
$2\max\{(1-p)^2,p^2\}$. Yet the full-support exchangeable law
$(\epsilon/2,(1-\epsilon)/2,(1-\epsilon)/2,\epsilon/2)$ gives
$\epsilon$ and $1-\epsilon$ for $0<\epsilon<1/2$.
The research proof gives an embedding in a single four-player game; it does
not assume all four roles have equal action counts.

\paragraph{Unknown risk attitude.}
For common ordered physical consequences $z_0<\cdots<z_L$, replace a row by
its upper-tail vector $(\Pr(r\ge z_\ell))_{\ell=1}^L$ under the fixed belief.
Equality of the downward convex hulls of these vectors is necessary and
sufficient for equal optimized values under every strictly increasing
expected-utility function on these consequences. Indeed utility increments
are arbitrary positive coefficients of the tails; continuity and separation
complete the argument. This uses common physical consequences and common
risk attitudes. It is not a characterization by role-specific ordinal
utility automorphisms.

\paragraph{Action-dependent opponents.}
The geometry extends to a disclosed response policy, with a different row
set. Let $c\in C_i$ be the participant's commitment before an unobserved
scenario $z$, and put $x_{ic}(z)=\mathbb E[W(r)\mid i,c,z]$. For finite
commitments, or compact commitments with continuous rows, equality of
$\operatorname{conv}\{x_{ic}\}+\R_-^Z$ still characterizes equality of
optimized values for every scenario law: compactness and the same separation
proof suffice. This does not replace the response-adjusted rows by the
original payoff table. For example, an opponent observing an announced
binary mixture $p$ can respond with success probabilities $p$ and $1-p^2$
in two scenarios. With receipt equal to success, the equal-prior optimum is
$5/8$ at $p=1/2$, while both pure commitments give $1/2$.
The response is a stipulated policy, not asserted to be a rational best response.
Nash or correlated-equilibrium behavior instead requires a specified equilibrium
correspondence and selection rule. Exact symmetry transports those sets;
it does not make every selected equilibrium role-indifferent.
